\section{An application example: IMU-Enabled Posture Monitor}

Paper \cite{petropoulos2020wearable} reports an interesting application for synchronized sensor nodes that is worth analyzing in this thesis: a monitoring system for sitting posture that uses wireless motion sensors attached to the user's back. Different sensors are used, and data coming from them is then fused together.

\paragraph{Implementation} The article focuses on three main points:
\begin{itemize}
    \item \textbf{Data Sensing} The required hardware needs to be small, comfortable to the user and compatible with all his possible movements. Of course, it needs to be powered via rechargeable batteries. All these requirements are the typical ones of IoT applications, and are supported by the new MEMS-based IMUs: small, low-power devices that provide all the necessary information.

    In this application, 6-degrees of freedom IMU sensors (namely, accelerometer and gyroscope) are installed on two prototypes modules that are attached to the upper and lower back of the user, parallel to their spine. The stream of these devices will then be used and elaborated to provide the spine misalignment devised into a pair of angles. Starting from this data, the system is capable of detecting spine inclination and lumbar placement in all the four directions. 

    Two main problems in using IMUs in this application have been assessed: first of all, the low-resolution, noise-corrupted and biased signals coming from these kind of sensors require a quite heavy processing stage to get quality data. Secondly, as it is only possible to detect angle changes when using accelerometer and gyroscope sensors, a calibration phase (with the user collaboration) is necessary before starting to harvest posture data.

    \item \textbf{Data Transmission} In this application, Bluetooth Low Energy has been considered the best solution. It has been increasingly used in wearable applications, it offers power-efficient operation and, mostly important, it is natively compatible with a lot of devices that could work as data aggregator, like smartphones.

    \item \textbf{Data Processing} As stated before, data coming from these IMU sensors requires some sort of refining to clean the signal and get the interesting angle information. Raw data is processed by means of a weighted moving average filter, then an Explicit Complementary Filter (ECF) algorithm is used to fuse all the information together and extract the two angles. All of this is performed on the smartphone of the user, that provides the required calculation capabilities (not available on the sensor nodes). 
\end{itemize}

\paragraph{User Interface} In order to use the posture monitor system, the user will first participate in a placement and calibration phase. Once it's done, the monitoring session starts, and the user is notified via vibration every time the angle threshold indicating a bad posture are violated. Once the session is complete, the user can analyze on the smartphone their posture behaviour, while the data is uploaded to a cloud-based application that can monitor the posture over time, maybe using Artificial Intelligence algorithms. 

\paragraph{Evaluation} To perform evaluation of the entire system, the two sensor boards have been attached on Pan/Tilt units. Then, different pre-defined couple of angles are simulated by the servo motors of the units, and then compared by the ones estimated by the IMUs. Repeated tests showed that the system is very reliable, showing errors that are distributed in a Gaussian shape and that are less than 0.1 degrees. The paper explains that observed errors are caused by the intrinsic error noise, and that the IMU calibration procedure is very useful to reduce biases, misalignment and gain errors. 

\paragraph{Conclusions} The article showed how data coming from low-cost sensor nodes, with limited computational capabilities and restrictive power constraint, can be quite precise when dealing with medical applications. Of course, there is the need of extensive post-processing work, usually performed by algorithms that run on more powerful devices, but the result is an user-friendly application that can be used by anyone.  